% UBT-AI-PROVENANCE-BEGIN
% schema: ubt-ai-provenance/v1
% tier: C_working
% ai_assistance: disclosed
% human_review: risk-based
% editorial_responsibility: Ing. David Jaroš
% policy: ../../../AI_PROVENANCE.md
% notice: Working material; exhaustive human review is not claimed.
% UBT-AI-PROVENANCE-END

\chapter{UBT za hodinu: motivace, architektura a současný stav}
\label{ch:overview}

\section{K čemu tato kapitola slouží}
Unified Biquaternion Theory (UBT) zkoumá, zda lze geometrii prostoročasu a
vnitřní strukturu polí zakódovat do jediného hlavního pole s hodnotami v
komplexifikovaných kvaternionech, místo aby byly zavedeny jako nesouvisející
základní objekty. Výchozí zápis současné teorie je
\[
 \Theta(q,\tau)\in\mathbb B=\mathbb C\otimes\mathbb H,
 \qquad \tau=t+i\psi .
\]
Zápis je kompaktní, ale skrývá několik logicky odlišných otázek:
co je definováno, co lze kinematicky reprezentovat, co vyplývá z akce a co je
pouze výzkumnou hypotézou. Učebnice tyto úrovně důsledně odděluje.

\section{Současný lokální geometrický řetězec}
Čistá lokální větev GR vychází z kovariantní derivace jediného pole,
\[
 E_\mu:=\mathcal N_0^{-1/2}D_\mu\Theta .
\]
Na reálném Lorentzově řezu píšeme
\[
 E_\mu=i e_\mu{}^0\mathbf1+e_\mu{}^k\mathbf e_k,
 \qquad e_\mu{}^a\in\mathbb R.
\]
Kvaternionovou konjugaci značíme $\sharp$. Symetrizovaný součin je centrální
a definuje metriku:
\[
 \boxed{
 \frac12\left(E_\mu^\sharp E_\nu+E_\nu^\sharp E_\mu\right)
 =g_{\mu\nu}\mathbf1 .}
\]
Minimální lokální konstrukce metriky tedy nevyžaduje projekci na pozorovatele,
průměrování přes kompaktní vlákno ani vkládací mapu. Nahrazuje tím starší
architektonický obraz dochovaný v historickém archivu.

Tatáž první derivace pole se chová jako tetráda. Pro nedegenerovanou tetrádu
má zobrazení $e_\mu{}^a\mapsto g_{\mu\nu}$ známý počet nezávislých složek
\[
 16-6=10,
\]
protože šest lokálních Lorentzových transformací ponechává deset složek
metriky beze změny. Tím se na kinematické úrovni řeší lokální problém hodnosti.

\section{Konexe, křivost a co je stále dynamické}
Lokální Lorentzova neboli bikvaternionová konexe vstupuje do $D_\mu$. Ve větvi
GR bez torze je kompatibilní spinovou konexí Leviho--Civitova spinová konexe
vypočtená z tetrády. Je-li torze nenulová, platí
\[
 \omega=\mathring\omega(e)+K(T),
\]
takže tetráda spolu s torzí určují metricky kompatibilní konexi. Křivost je
obsažena v komutátoru kovariantních derivací, nikoli v algebraickém komutátoru
vektorů tetrády:
\[
 \mathcal R_{\mu\nu}:=[D_\mu,D_\nu].
\]

Současná práce na GR uzavírá důležité otázky lokální reprezentovatelnosti,
dosud však neodvozuje úplnou Einsteinovu dynamiku z původní hlavní akce UBT.
Výběr akce, fyzikální omezení torze, úplné propojení s poruchovou teorií a
globální pokračování proto zůstávají otevřenými problémy s výslovně vedeným
stavem. Směrodatné znění je uvedeno v souborech
\texttt{STATUS.md} a \texttt{WHAT\_IS\_PROVED.md}.

\section{Komplexní čas a theta/spektrální větev}
Kanonický zápis zachovává vztah
\[
 \tau=t+i\psi,
\]
přičemž $\psi$ se v aktivní větvi komplexního času interpretuje jako vnitřní
fázová souřadnice, nikoli jako druhý makroskopický čas. Redukované theta modely
pak přirozeně obsahují členy tvaru
\[
 S(t,\psi)=\sum_n a_n e^{-\pi\psi n^2}e^{i\pi t n^2}.
\]
Jediný výraz tak připouští několik standardních reprezentací: Fourierovu v
$t$, Laplaceovu či Mellinovu v $\psi$, Jacobiho--Poissonovu dualitu a v
kontrolovaných kvadratických příkladech také vztah mezi součty módů a součty
dráhových či topologických tříd. Tyto klasické nástroje jsou podrobně rozvedeny
v dalších kapitolách, protože propojují několik jinak oddělených směrů výzkumu
UBT.

\section{Kalibrační a částicové sektory}
Algebra $\mathbb C\otimes\mathbb H\cong M_2(\mathbb C)$ přirozeně podporuje
levé i pravé působení, spinorové Lorentzovy transformace a konstrukce vnitřních
symetrií. Některá tvrzení o kalibračním sektoru jsou na uvedených úrovních
dokázána nebo podmíněně uzavřena; jiná zůstávají otevřená, včetně částí
hypernáboje, výběru generací a dynamického narušení symetrie. Učebnice proto z
pouhé existence bikvaternionové algebry nevyvozuje úplné odvození standardního
modelu.

\section{Konstanta jemné struktury: současné bezpečné tvrzení}
Program alfa obsahuje reprodukovatelný strukturální mechanismus, který vybírá
celé číslo $137$ za uvedených předpokladů, včetně výsledku o stabilitě vůči
prvočíselným kandidátům. Celý výzkumný směr však v současnosti nepředstavuje
odvození fyzikální konstanty jemné struktury z prvních principů. Zvlášť
důležité jsou dvě překážky:
\begin{enumerate}
 \item efektivní koeficient požadovaný potenciálem vinutí dosud nebyl
       jednoznačně odvozen z akce UBT (mezera G137-B);
 \item auditovaná násobnost smyček a násobnost sektoru se zkroucením dosud
       nebyly ztotožněny dokázaným propojením.
\end{enumerate}
Proto zde nelze ani označení „odvození alfa bez přizpůsobování parametrů“, ani číselnou hodnotu
$137.036\ldots$ prezentovat jako uzavřenou předpověď UBT.
Kapitola~\ref{ch:alpha} odvozuje podmíněně platnou matematiku a v každém kroku
uvádí zbývající mezery.

\section{Jak číst značky stavu}
Repozitář používá dva nezávislé systémy evidence, které se nesmějí zaměňovat.

Úrovně vědeckých tvrzení rozlišují standardní identity, dokázaná lemmata UBT,
podmíněné výsledky, motivované domněnky, pozorování a otevřené mezery. Úrovně
provenience AI naproti tomu zaznamenávají redakční stav a míru kontroly souborů.
Definuje je \texttt{PROVENANCE\_TIERS.yaml}: úroveň A potvrzuje autor, úroveň B
je strojově ověřená, úroveň C je pracovní a úroveň D historická. Tato učebnice
je pracovním materiálem úrovně C. Jejím účelem je vysvětlit a propojit
matematiku, nikoli bez dokladů povýšit některé tvrzení na úroveň A.

\section{Pravidlo autoritativního zdroje pro učebnici}
Zdrojová mapa v \texttt{docs/textbook/SOURCE\_MAP.md} zaznamenává, odkud každá
kapitola přebírá aktuální stav tvrzení. Historický materiál v
\texttt{ARCHIVE/} může být cenný pro pochopení vývoje určité myšlenky, nikdy se
však bez výslovného označení nevkládá příkazem
\texttt{\textbackslash input} jako součást současné teorie.
